\chapter{Berechnungstheoretische Vorteile bestimmter Graphklassen zur Bestimmung bedingter Unabhängigkeit}
\label{ch:graphclasses}

Die folgenden Abschnitte orientieren sich an \cite{handbook__of_formal_argumentation__vol__3:2024}.

\section{Azyklische AFs}
\label{sec:acyc-afs}

Für azyklische AFs fällt jedes existierende \emph{stable}-Labelling mit dem eindeutig bestimmten \emph{grounded}-Labelling zusammen. 
Folglich besitzt ein azyklisches AF höchstens ein \emph{stable}-Labelling. 
Dadurch ist die bedingte Unabhängigkeit bezüglich der \emph{stable}-Semantik für beliebige paarweise disjunkte Argumentmengen 
stets erfüllt.

\begin{Proposition}
\label{prop:acyc-af-ci-st-top}
    Sei F=(A,R) ein azyklisches AF. Dann ist
    \(\textit{Ind}_{\mathrm{st}}(F) \) immer wahr.
\end{Proposition}

\begin{proof}
    \label{pr:acyc-af-ci-st-top-true}
    Seien \( X,Y,Z \subseteq A\) paarweise disjunkt. Besitzt \(F\) kein
    \emph{stable}-Labelling, so ist die universell quantifizierte
    Unabhängigkeitsbedingung vakuos erfüllt.

    Andernfalls existiert aufgrund der Azyklizität genau ein
    \emph{stable}-Labelling \(\lambda\), das mit dem
    \emph{grounded}-Labelling übereinstimmt. Für beliebige
    \(\lambda_1,\lambda_2\in\Lambda_{\mathrm{st}}(F)\) gilt daher
    \(\lambda_1=\lambda_2=\lambda\). Mit \(\lambda_3:=\lambda\) können
    somit die Belegungen von \(X\), \(Y\) und \(Z\) wie in der Definition
    der bedingten Unabhängigkeit gefordert kombiniert werden. Daher gilt
    \(X\bot_{\mathrm{st}}Y\mid Z\) und folglich
    \(\operatorname{Ind}_{\mathrm{st}}(F)\).
\end{proof}

\section{Geradezykelfreie AFs}
\label{sec:even-cycle-free-afs}

Für AFs ohne gerade Zykel ist allein das \emph{grounded}-Labelling ein Kandidat, ein \emph{stable}-Labelling zu sein. 
Damit existiert ebenfalls maximal ein \emph{stable}-Labelling. 

\begin{Proposition}
\label{prop:even-cycle-free-af-ci-st-top}
    Sei F=(A,R) ein AF ohne gerade Zyklen. Dann ist
    \(\textit{Ind}_{\mathrm{st}}(F) \) immer wahr.
\end{Proposition}

\begin{proof}
    Analog zu \ref{pr:acyc-af-ci-st-top-true}.
\end{proof}


\section{Ungeradezykelfreie AFs}
\label{sec:odd-cycle-free-afs}

Für die Berechnung von \emph{stable}-Labellings gibt es in dieser Graphklasse keine berechnungstheoretischen Vorteile. 


\section{Bipartite AFs}
\label{sec:bi-afs}

Als Spezialfall eines ungeradezykelfreien AFs gibt es für die Berechnung von \emph{stable}-Labellings in bipartiten AFs keine
berechnungstheoretischen Vorteile. 


\section{Symmetrische AFs}
\label{sec:sym-afs}

Beobachtung: Je mehr Angriffe im AF, desto häufiger nicht bedingt unabhängig.

